

% Page Layout
\usepackage[tmargin=2cm,rmargin=1in,lmargin=1in,margin=0.85in,bmargin=2cm,footskip=0.0in]{geometry}
% Controls page margins, text width/height, footskip, etc.

\setlength{\parindent}{0pt}      % no paragraph indent
\setlength{\parskip}{0.4em}      % small, consistent gap between paragraphs


% 📚 Math & Symbols
\usepackage{amsmath,amsfonts,amsthm,amssymb,mathtools,bbm,bm}
% Core AMS math tools (theorems, align, fonts, symbols).

% bbm – Provides blackboard bold \mathbbm{1} indicator symbol.
% bm – Bold math symbols (\bm{x}).
% newpxmath – Times-like math font compatible with Palatino/PX fonts.
% xfrac – Inline slanted fractions using \sfrac{a}{b}.
% cancel – Strikeout terms in equations with \cancel{}.
\usepackage[varbb]{newpxmath}

\usepackage{xfrac}
\usepackage[makeroom]{cancel}
\usepackage{mathtools}



% 🔗 Navigation & References
\usepackage{bookmark} % Improved handling of PDF bookmarks.
\usepackage{hyperref,theoremref} 
% Allows references like “Theorem 1.2”.
\usepackage{nameref} 
% Reference by section name (\nameref{section}).



% 📦 Layout, Lists, Tables
\usepackage{enumitem} % Customize list spacing and labels
\usepackage{multicol,array} 
% multicol – Multiple column layout inside the page.
% array – Enhanced tabular column features.



% 🎨 Graphics & Coloring
\usepackage{xcolor} % Colored text and drawing.
\usepackage{tikz-cd} % Commutative diagrams (TikZ-based).
\usepackage{transparent} % Adjust text/image transparency.
% \usepackage[most]{tcolorbox}


% 📑 Frames, Boxes, and Visual Elements
\usepackage[most,many,breakable]{tcolorbox}
% Colored/fancy boxes, breakable boxes, theorems.
\usepackage{varwidth} % Boxes with variable width (auto-fit content)
% <<< PUT IT RIGHT HERE >>>
\tcbset{fontupper=\normalsize}


% 🔧 Utilities & Tools
\usepackage{etoolbox}
% Logic programming tools (patching commands, conditionals).
\usepackage{xifthen} 
% Conditional statements (\ifthenelse) with logic.
\usepackage{import}
% Import figures or files from subdirectories.
\usepackage{pdfpages}
% Insert external PDF pages (\includepdf).
\usepackage{comment} 
% Multi-line comment environment.
\usepackage[ruled,vlined,linesnumbered]{algorithm2e}
% Nice-looking algorithm environment (algorithm2e).
\usepackage{natbib}
% Author–year and numeric citation management (needed for JFE, JFQA styles).

\usepackage{tikz}
\usetikzlibrary{positioning}
\usepackage{tikzsymbols}
\renewcommand\qedsymbol{$\Laughey$}

% the package for the table
% \usepackage{booktabs}
% \usepackage{threeparttable}
\usepackage{multirow}
% \usepackage{adjustbox}
\usepackage{tabularx}
\newcolumntype{Y}{>{\raggedleft\arraybackslash}X}


% minted is a LaTeX package for beautifully formatting source code with syntax highlighting.
% \usepackage{minted}
% \definecolor{LightGray}{gray}{0.95}

% \hypersetup{
% 	pdftitle={Notes},
% 	colorlinks=true, linkcolor=blue,
% 	bookmarksnumbered=true,
% 	bookmarksopen=true
% }

% \hypersetup{
%     pdftitle={PhD Research Notes},
%     colorlinks=true,
%     linkcolor=blue!50!black,
%     citecolor=blue!50!black,
%     urlcolor=blue!70!black,
%     bookmarksnumbered=true,
%     bookmarksopen=true
% }

% \hypersetup{
%     pdftitle={PhD Research Notes},
%     colorlinks=true,
%     linkcolor=blue,
%     citecolor=blue,
%     urlcolor=blue,
%     bookmarksnumbered=true,
%     bookmarksopen=true
% }

% Original hyper set up
\hypersetup{
	pdftitle={Assignment},
	colorlinks=true, linkcolor=doc!90,
	bookmarksnumbered=true,
	bookmarksopen=true
}

% Nice tables
\usepackage{booktabs}        % \toprule \midrule \bottomrule
\usepackage{threeparttable}  % notes under tables
\usepackage{siunitx}         % numeric alignment
\usepackage{adjustbox}       % resize table to fit width
\usepackage{caption}         % caption styling (optional)

% draw figure in latex
\usepackage{tikz}
\usetikzlibrary{arrows.meta, positioning, shapes.geometric}


% siunitx setup: adjust to your taste
\sisetup{
  detect-all,
  input-symbols   = (),
  table-number-alignment = center,
  group-separator = {,},
  group-minimum-digits = 4,
  table-space-text-post = ***,
}

% slightly tighter / nicer table spacing
\setlength{\tabcolsep}{6pt}
\renewcommand{\arraystretch}{1.15}


% %%%%%%%%%%%%%%%%%%%%%%%%%%%%%%
% % SELF MADE COLORS
% %%%%%%%%%%%%%%%%%%%%%%%%%%%%%%



\definecolor{myg}{RGB}{56, 140, 70}
\definecolor{myb}{RGB}{45, 111, 177}
% \definecolor{myr}{RGB}{199, 68, 64}
\definecolor{mytheorembg}{HTML}{F2F2F9}
\definecolor{mytheoremfr}{HTML}{00007B}
% \definecolor{mylenmabg}{HTML}{FFFAF8}
% \definecolor{mylenmafr}{HTML}{983b0f}
% \definecolor{mypropbg}{HTML}{f2fbfc}
% \definecolor{mypropfr}{HTML}{191971}
\definecolor{myexamplebg}{HTML}{F2FBF8}
\definecolor{myexamplefr}{HTML}{88D6D1}
\definecolor{myexampleti}{HTML}{2A7F7F}
% \definecolor{mydefinitbg}{HTML}{E5E5FF}
% \definecolor{mydefinitfr}{HTML}{3F3FA3}
% \definecolor{notesgreen}{RGB}{0,162,0}
\definecolor{myp}{RGB}{197, 92, 212}
\definecolor{mygr}{HTML}{2C3338}
\definecolor{myred}{RGB}{127,0,0}
\definecolor{myyellow}{RGB}{169,121,69}
% \definecolor{myexercisebg}{HTML}{F2FBF8}
% \definecolor{myexercisefg}{HTML}{88D6D1}

\definecolor{softredbg}{RGB}{255,235,238}   % very light, soft red background
\definecolor{softredline}{RGB}{198,40,40}   % deep soft red for the border/title
\definecolor{softbluebg}{RGB}{232,245,253}   % very light soft-blue background
\definecolor{softblueline}{RGB}{30,136,229}  % soft but strong blue (border/title)
\definecolor{softgreybg}{RGB}{245,245,245}   % very light grey background
\definecolor{softgreyline}{RGB}{150,150,150} % medium grey border/title


% =====================================================
%  COLOR PALETTE FOR READING NOTES
% =====================================================
\definecolor{refabg}{RGB}{232,245,233}   % soft blue     (Abstract)
\definecolor{refaline}{RGB}{56,142,60}

\definecolor{refmbg}{RGB}{232,245,253}   % soft green    (Method)
\definecolor{refmline}{RGB}{30,136,229}

\definecolor{refcbg}{RGB}{255,235,238}   % soft red   (Conclusion)
\definecolor{refcline}{RGB}{198,40,40}

\definecolor{reflbg}{RGB}{255,249,230}   % soft yellow   (Background/Literature)
\definecolor{reflline}{RGB}{251,192,45}

\definecolor{refnote_bg}{RGB}{245,245,245}   % soft grey (Your comments)
\definecolor{refnote_line}{RGB}{120,120,120}

\tcbuselibrary{skins,breakable}

% =====================================================
% 1. ABSTRACT BOX
% =====================================================
\newtcolorbox{refabstract}[1][]{%
  enhanced, breakable,
  colback=refabg,
  colframe=refaline,
  boxrule=0.8pt,
  borderline west={2pt}{0pt}{refaline},
  sharp corners,
  title=\textbf{Abstract},
  coltitle=refaline,
  fonttitle=\bfseries\sffamily,
  detach title,
  before upper=\tcbtitle\par\smallskip,
  #1
}

% =====================================================
% 2. METHODOLOGY BOX
% =====================================================
\newtcolorbox{refmethod}[1][]{%
  enhanced, breakable,
  colback=refmbg,
  colframe=refmline,
  boxrule=0.8pt,
  borderline west={2pt}{0pt}{refmline},
  sharp corners,
  title=\textbf{Methodology},
  coltitle=refmline,
  fonttitle=\bfseries\sffamily,
  detach title,
  before upper=\tcbtitle\par\smallskip,
  #1
}

% =====================================================
% 3. KEY FINDINGS / CONCLUSION BOX
% =====================================================
\newtcolorbox{refconclusion}[1][]{%
  enhanced, breakable,
  colback=refcbg,
  colframe=refcline,
  boxrule=0.8pt,
  borderline west={2pt}{0pt}{refcline},
  sharp corners,
  title=\textbf{Key Findings / Conclusion},
  coltitle=refcline,
  fonttitle=\bfseries\sffamily,
  detach title,
  before upper=\tcbtitle\par\smallskip,
  #1
}

% =====================================================
% 4. BACKGROUND / LITERATURE / MOTIVATION BOX
% =====================================================
\newtcolorbox{reflitnote}[1][]{%
  enhanced, breakable,
  colback=reflbg,
  colframe=reflline,
  boxrule=0.8pt,
  borderline west={2pt}{0pt}{reflline},
  sharp corners,
  title=\textbf{Background / Literature},
  coltitle=reflline,
  fonttitle=\bfseries\sffamily,
  detach title,
  before upper=\tcbtitle\par\smallskip,
  #1
}

% =====================================================
% 5. YOUR OWN COMMENTS / CRITIQUE BOX
% =====================================================
\newtcolorbox{refcomment}[1][]{%
  enhanced, breakable,
  colback=refnote_bg,
  colframe=refnote_line,
  boxrule=0.8pt,
  borderline west={2pt}{0pt}{refnote_line},
  sharp corners,
  title=\textbf{My Comments},
  coltitle=refnote_line,
  fonttitle=\bfseries\sffamily,
  detach title,
  before upper=\tcbtitle\par\smallskip,
  #1
}


\newtcbox{\mybox}[1][red]{on line,
arc=0pt,outer arc=0pt,colback=#1!10!white,colframe=#1!50!black,
boxsep=0pt,left=1pt,right=1pt,top=2pt,bottom=2pt,
boxrule=0pt,bottomrule=1pt,toprule=1pt}
\newtcbox{\xmybox}[1][red]{on line,
arc=7pt,colback=#1!10!white,colframe=#1!50!black,
before upper={\rule[-3pt]{0pt}{10pt}},boxrule=1pt,
boxsep=0pt,left=6pt,right=6pt,top=2pt,bottom=2pt}


%================================
% EXAMPLE BOX
%================================

\newtcbtheorem[number within=section]{Example}{Example}
{%
	colback = myexamplebg
	,breakable
	,colframe = myexamplefr
	,coltitle = myexampleti
	,boxrule = 1pt
	,sharp corners
	,detach title
	,before upper=\tcbtitle\par\smallskip
	,fonttitle = \bfseries
	,description font = \mdseries
	,separator sign none
	,description delimiters parenthesis
}
{ex}

\newtcbtheorem[number within=chapter]{example}{Example}
{%
	colback = myexamplebg
	,breakable
	,colframe = myexamplefr
	,coltitle = myexampleti
	,boxrule = 1pt
	,sharp corners
	,detach title
	,before upper=\tcbtitle\par\smallskip
	,fonttitle = \bfseries
	,description font = \mdseries
	,separator sign none
	,description delimiters parenthesis
}
{ex}


%================================
% Question BOX
%================================

\makeatletter
\newtcbtheorem{question}{Question}{enhanced,
	breakable,
	colback=white,
	colframe=myb!80!black,
	attach boxed title to top left={yshift*=-\tcboxedtitleheight},
	fonttitle=\bfseries,
	title={#2},
	boxed title size=title,
	boxed title style={%
			sharp corners,
			rounded corners=northwest,
			colback=tcbcolframe,
			boxrule=0pt,
		},
	underlay boxed title={%
			\path[fill=tcbcolframe] (title.south west)--(title.south east)
			to[out=0, in=180] ([xshift=5mm]title.east)--
			(title.center-|frame.east)
			[rounded corners=\kvtcb@arc] |-
			(frame.north) -| cycle;
		},
	#1
}{def}
\makeatother



%================================
% SOLUTION BOX
%================================

\makeatletter
\newtcolorbox{solution}{enhanced,
	breakable,
	colback=white,
	colframe=myg!80!black,
	attach boxed title to top left={yshift*=-\tcboxedtitleheight},
	title=Solution,
	boxed title size=title,
	boxed title style={%
			sharp corners,
			rounded corners=northwest,
			colback=tcbcolframe,
			boxrule=0pt,
		},
	underlay boxed title={%
			\path[fill=tcbcolframe] (title.south west)--(title.south east)
			to[out=0, in=180] ([xshift=5mm]title.east)--
			(title.center-|frame.east)
			[rounded corners=\kvtcb@arc] |-
			(frame.north) -| cycle;
		},
}
\makeatother

% %================================
% % Question BOX
% %================================

% \makeatletter
% \newtcbtheorem{qstion}{Question}{enhanced,
% 	breakable,
% 	colback=white,
% 	colframe=mygr,
% 	attach boxed title to top left={yshift*=-\tcboxedtitleheight},
% 	fonttitle=\bfseries,
% 	title={#2},
% 	boxed title size=title,
% 	boxed title style={%
% 			sharp corners,
% 			rounded corners=northwest,
% 			colback=tcbcolframe,
% 			boxrule=0pt,
% 		},
% 	underlay boxed title={%
% 			\path[fill=tcbcolframe] (title.south west)--(title.south east)
% 			to[out=0, in=180] ([xshift=5mm]title.east)--
% 			(title.center-|frame.east)
% 			[rounded corners=\kvtcb@arc] |-
% 			(frame.north) -| cycle;
% 		},
% 	#1
% }{def}
% \makeatother

% \newtcbtheorem[number within=chapter]{wconc}{Wrong Concept}{
% 	breakable,
% 	enhanced,
% 	colback=white,
% 	colframe=myr,
% 	arc=0pt,
% 	outer arc=0pt,
% 	fonttitle=\bfseries\sffamily\large,
% 	colbacktitle=myr,
% 	attach boxed title to top left={},
% 	boxed title style={
% 			enhanced,
% 			skin=enhancedfirst jigsaw,
% 			arc=3pt,
% 			bottom=0pt,
% 			interior style={fill=myr}
% 		},
% 	#1
% }{def}

% %%%%%%%%%%%%%%%%%%%%%%%%%%%%
% % TCOLORBOX SETUPS
% %%%%%%%%%%%%%%%%%%%%%%%%%%%%

\setlength{\parindent}{0pt}

%================================
% Brainteasers
%================================

\tcbuselibrary{theorems,skins,hooks}
\newtcbtheorem[number within=section]{brainteasers}{Brainteasers}
{%
	enhanced
	,breakable
	,colback = myg!10
	,frame hidden
	,boxrule = 0sp
	,borderline west = {2pt}{0pt}{myg}
	,sharp corners
	,detach title
	,before upper = \tcbtitle\par\smallskip
	,coltitle = myg!85!black
	,fonttitle = \bfseries\sffamily
	,description font = \mdseries
	,separator sign none
	,segmentation style={solid, myg!85!black}
}
{th}


% %================================
% % CLAIM
% %================================

% \tcbuselibrary{theorems,skins,hooks}
% \newtcbtheorem[number within=section]{claim}{Literature}
% {%
% 	enhanced
% 	,breakable
% 	,colback = myg!10
% 	,frame hidden
% 	,boxrule = 0sp
% 	,borderline west = {2pt}{0pt}{myg}
% 	,sharp corners
% 	,detach title
% 	,before upper = \tcbtitle\par\smallskip
% 	,coltitle = myg!85!black
% 	,fonttitle = \bfseries\sffamily
% 	,description font = \mdseries
% 	,separator sign none
% 	,segmentation style={solid, myg!85!black}
% }
% {th}

% %================================
% % SOFT BLUE COLORS
% %================================
% \definecolor{softbluebg}{RGB}{232,245,253}   % light blue background
% \definecolor{softblueline}{RGB}{30,136,229}  % soft strong blue border/title

%================================
% THEOREM BOX
%================================

\tcbuselibrary{theorems,skins,hooks}
\newtcbtheorem[number within=section]{Theorem}{Theorem}
{%
	enhanced,
	breakable,
	colback = mytheorembg,
	frame hidden,
	boxrule = 0sp,
	borderline west = {2pt}{0pt}{mytheoremfr},
	sharp corners,
	detach title,
	before upper = \tcbtitle\par\smallskip,
	coltitle = mytheoremfr,
	fonttitle = \bfseries\sffamily,
	description font = \mdseries,
	separator sign none,
	segmentation style={solid, mytheoremfr},
}
{th}

% \tcbuselibrary{theorems,skins,hooks}
% \newtcbtheorem[number within=chapter]{theorem}{Theorem}
% {%
% 	enhanced,
% 	breakable,
% 	colback = mytheorembg,
% 	frame hidden,
% 	boxrule = 0sp,
% 	borderline west = {2pt}{0pt}{mytheoremfr},
% 	sharp corners,
% 	detach title,
% 	before upper = \tcbtitle\par\smallskip,
% 	coltitle = mytheoremfr,
% 	fonttitle = \bfseries\sffamily,
% 	description font = \mdseries,
% 	separator sign none,
% 	segmentation style={solid, mytheoremfr},
% }
% {th}


\tcbuselibrary{theorems,skins,hooks}
\newtcolorbox{Theoremcon}
{%
	enhanced
	,breakable
	,colback = mytheorembg
	,frame hidden
	,boxrule = 0sp
	,borderline west = {2pt}{0pt}{mytheoremfr}
	,sharp corners
	,description font = \mdseries
	,separator sign none
}


\tcbuselibrary{theorems,skins,hooks}

\newtcbtheorem[number within=section]{problem}{Problem}
{%
    enhanced,
    breakable,
    colback = mytheorembg,
    frame hidden,
    boxrule = 0sp,
    borderline west = {2pt}{0pt}{mytheoremfr},
    sharp corners,
    detach title,
    before upper = \tcbtitle\par\smallskip,
    coltitle = mytheoremfr,
    fonttitle = \bfseries\sffamily,
    description font = \mdseries,
    separator sign none,
    segmentation style={solid, mytheoremfr},
}
{prob}



% %================================
% % THEOREM BOX (Methodology)
% %================================
% \tcbuselibrary{theorems,skins,hooks}
% \newtcbtheorem[number within=section]{Theorem}{Methodology}
% {%
%     enhanced,
%     breakable,
%     colback = softbluebg,                     % background
%     frame hidden,
%     boxrule = 0sp,
%     borderline west = {2pt}{0pt}{softblueline}, % left stripe
%     sharp corners,
%     detach title,
%     before upper = \tcbtitle\par\smallskip,
%     coltitle = softblueline,                  % title color
%     fonttitle = \bfseries\sffamily,
%     description font = \mdseries,
%     separator sign none,
%     segmentation style={solid, softblueline}, % internal segmentation line
% }
% {th}

% %================================
% % THEOREM BOX (Literature Review)
% %================================
% \tcbuselibrary{theorems,skins,hooks}
% \newtcbtheorem[number within=chapter]{theorem}{Literature}
% {%
%     enhanced,
%     breakable,
%     colback = softbluebg,
%     frame hidden,
%     boxrule = 0sp,
%     borderline west = {2pt}{0pt}{softblueline},
%     sharp corners,
%     detach title,
%     before upper = \tcbtitle\par\smallskip,
%     coltitle = softblueline,
%     fonttitle = \bfseries\sffamily,
%     description font = \mdseries,
%     separator sign none,
%     segmentation style={solid, softblueline},
% }
% {th}

% %================================
% % THEOREM BOX (Theoremcon)
% %================================
% \tcbuselibrary{theorems,skins,hooks}
% \newtcolorbox{Theoremcon}
% {%
%     enhanced,
%     breakable,
%     colback = softbluebg,
%     frame hidden,
%     boxrule = 0sp,
%     borderline west = {2pt}{0pt}{softblueline},
%     sharp corners,
%     description font = \mdseries,
%     separator sign none
% }


% %================================
% % Corollery (Conclusion)
% %================================
% \tcbuselibrary{theorems,skins,hooks}
% \newtcbtheorem[number within=section]{Corollary}{Conclusion}
% {%
%     enhanced,
%     breakable,
%     colback = softredbg,                       % background
%     frame hidden,
%     boxrule = 0sp,
%     borderline west = {2pt}{0pt}{softredline}, % left stripe
%     sharp corners,
%     detach title,
%     before upper = \tcbtitle\par\smallskip,
%     coltitle = softredline,                    % title text color
%     fonttitle = \bfseries\sffamily,
%     description font = \mdseries,
%     separator sign none,
%     segmentation style={solid, softredline},   % internal segment line
% }
% {th}

% % \tcbuselibrary{theorems,skins,hooks}
% % \newtcbtheorem[number within=chapter]{corollary}{Conclusion}
% % {%
% % 	enhanced
% % 	,breakable
% % 	,colback = myp!10
% % 	,frame hidden
% % 	,boxrule = 0sp
% % 	,borderline west = {2pt}{0pt}{myp!85!black}
% % 	,sharp corners
% % 	,detach title
% % 	,before upper = \tcbtitle\par\smallskip
% % 	,coltitle = myp!85!black
% % 	,fonttitle = \bfseries\sffamily
% % 	,description font = \mdseries
% % 	,separator sign none
% % 	,segmentation style={solid, myp!85!black}
% % }
% % {th}



% % %================================
% % % NOTE BOX
% % %================================

% % \usetikzlibrary{arrows,calc,shadows.blur}
% % \tcbuselibrary{skins}
% % \newtcolorbox{note}[1][]{%
% % 	enhanced jigsaw,
% % 	colback=gray!20!white,%
% % 	colframe=gray!80!black,
% % 	size=small,
% % 	boxrule=1pt,
% % 	title=\textbf{Note:-},
% % 	halign title=flush center,
% % 	coltitle=black,
% % 	breakable,
% % 	drop shadow=black!50!white,
% % 	attach boxed title to top left={xshift=1cm,yshift=-\tcboxedtitleheight/2,yshifttext=-\tcboxedtitleheight/2},
% % 	minipage boxed title=1.5cm,
% % 	boxed title style={%
% % 			colback=white,
% % 			size=fbox,
% % 			boxrule=1pt,
% % 			boxsep=2pt,
% % 			underlay={%
% % 					\coordinate (dotA) at ($(interior.west) + (-0.5pt,0)$);
% % 					\coordinate (dotB) at ($(interior.east) + (0.5pt,0)$);
% % 					\begin{scope}
% % 						\clip (interior.north west) rectangle ([xshift=3ex]interior.east);
% % 						\filldraw [white, blur shadow={shadow opacity=60, shadow yshift=-.75ex}, rounded corners=2pt] (interior.north west) rectangle (interior.south east);
% % 					\end{scope}
% % 					\begin{scope}[gray!80!black]
% % 						\fill (dotA) circle (2pt);
% % 						\fill (dotB) circle (2pt);
% % 					\end{scope}
% % 				},
% % 		},
% % 	#1,
% % }


% %================================
% % LENMA (Abstract)
% %================================

% \tcbuselibrary{theorems,skins,hooks}
% \newtcbtheorem[number within=section]{Lenma}{Abstract}
% {%
% 	enhanced,
% 	breakable,
% 	colback = mylenmabg,
% 	frame hidden,
% 	boxrule = 0sp,
% 	borderline west = {2pt}{0pt}{mylenmafr},
% 	sharp corners,
% 	detach title,
% 	before upper = \tcbtitle\par\smallskip,
% 	coltitle = mylenmafr,
% 	fonttitle = \bfseries\sffamily,
% 	description font = \mdseries,
% 	separator sign none,
% 	segmentation style={solid, mylenmafr},
% }
% {th}

% % \tcbuselibrary{theorems,skins,hooks}
% % \newtcbtheorem[number within=chapter]{lenma}{Lenma}
% % {%
% % 	enhanced,
% % 	breakable,
% % 	colback = mylenmabg,
% % 	frame hidden,
% % 	boxrule = 0sp,
% % 	borderline west = {2pt}{0pt}{mylenmafr},
% % 	sharp corners,
% % 	detach title,
% % 	before upper = \tcbtitle\par\smallskip,
% % 	coltitle = mylenmafr,
% % 	fonttitle = \bfseries\sffamily,
% % 	description font = \mdseries,
% % 	separator sign none,
% % 	segmentation style={solid, mylenmafr},
% % }
% % {th}


% %================================
% % NOTE BOX (Soft Blue Theme)
% %================================

% \usetikzlibrary{arrows,calc,shadows.blur}
% \tcbuselibrary{skins}

% \newtcolorbox{note}[1][]{%
%     enhanced jigsaw,
%     colback = gray!20!white,               % background
%     colframe = gray!80!black,            % frame color
%     size = small,
%     boxrule = 1pt,
%     title = \textbf{Note:},
%     halign title = flush center,
%     coltitle = black,            % title color
%     breakable,
%     drop shadow = black!50!white, % soft blue shadow
%     attach boxed title to top left = {
%         xshift = 1cm,
%         yshift = -\tcboxedtitleheight/2,
%         yshifttext = -\tcboxedtitleheight/2
%     },
%     minipage boxed title = 1.5cm,
%     boxed title style = {
%         colback = white,
%         size = fbox,
%         boxrule = 1pt,
%         boxsep = 2pt,
%         underlay = {
%             \coordinate (dotA) at ($(interior.west) + (-0.5pt,0)$);
%             \coordinate (dotB) at ($(interior.east) + (0.5pt,0)$);
%             \begin{scope}
%                 \clip (interior.north west) rectangle ([xshift=3ex]interior.east);
%                 \filldraw [white, blur shadow={shadow opacity=60, shadow yshift=-.75ex}, rounded corners=2pt] (interior.north west) rectangle (interior.south east);
%             \end{scope}
%             \begin{scope}[gray!80!black]
%                 \fill (dotA) circle (2pt);
%                 \fill (dotB) circle (2pt);
%             \end{scope}
%         },
%     },
%     #1,
% }




% %%%%%%%%%%%%%%%%%%%%%%%%%%%%%%
% % SELF MADE COMMANDS
% %%%%%%%%%%%%%%%%%%%%%%%%%%%%%%


\newcommand{\thm}[2]{\begin{Theorem}{#1}{}#2\end{Theorem}}
% \newcommand{\cor}[2]{\begin{Corollary}{#1}{}#2\end{Corollary}}
% % \newcommand{\prope}[2]{\begin{Property}{#1}{}#2\end{Property}}
% \newcommand{\mlenma}[2]{\begin{Lenma}{#1}{}#2\end{Lenma}}
% % \newcommand{\mprop}[2]{\begin{Prop}{#1}{}#2\end{Prop}}
% \newcommand{\clm}[3]{\begin{claim}{#1}{#2}#3\end{claim}}
% % \newcommand{\wc}[2]{\begin{wconc}{#1}{}\setlength{\parindent}{1cm}#2\end{wconc}}
% % \newcommand{\thmcon}[1]{\begin{Theoremcon}{#1}\end{Theoremcon}}
\newcommand{\ex}[2]{\begin{Example}{#1}{}#2\end{Example}}
% % \newcommand{\dfn}[2]{\begin{Definition}[colbacktitle=red!75!black]{#1}{}#2\end{Definition}}
% % \newcommand{\dfnc}[2]{\begin{definition}[colbacktitle=red!75!black]{#1}{}#2\end{definition}}
\newcommand{\qs}[2]{\begin{question}{#1}{}#2\end{question}}
\newcommand{\pf}[2]{\begin{myproof}[#1]#2\end{myproof}}
% \newcommand{\nt}[1]{\begin{note}#1\end{note}}

\newcommand{\qprob}[2]{\begin{problem}{#1}{}#2\end{problem}}
\newcommand{\bts}[2]{\begin{brainteasers}{#1}{}#2\end{brainteasers}}
% \newcommand{\bts}[3]{\begin{brainteasers}{#1}{#2}#3\end{brainteasers}}

% \newcommand{\sol}{\setlength{\parindent}{0cm}\textbf{\textit{Solution:}}\setlength{\parindent}{1cm} }
% \newcommand{\solve}[1]{\setlength{\parindent}{0cm}\textbf{\textit{Solution: }}\setlength{\parindent}{1cm}#1 \Qed}

% \newenvironment{sol}{
%     \par\setlength{\parindent}{0pt}%
%     \textbf{\textit{Solution:}}\par
% }{%
%     \par\setlength{\parindent}{15pt}%
% }

% \newenvironment{sol}{
%     \par\noindent\textbf{\textit{Solution:}}\par
%     \setlength{\parindent}{30pt}
% }{%
%     \par\setlength{\parindent}{15pt}
%     \qedsymbol
% }

\newcommand{\sol}{\setlength{\parindent}{0cm}\textbf{\textit{Solution:}}\setlength{\parindent}{1cm} }
\newcommand{\solve}[1]{\setlength{\parindent}{0cm}\textbf{\textit{Solution: }}\setlength{\parindent}{1cm}#1 \Qed}

% \newenvironment{sol}{
%     \begin{trivlist}
%     \item \textbf{\textit{Solution:}}
%     \setlength{\parindent}{30pt}
%     \begin{minipage}{\linewidth}
% }{
%     \vspace{0.2em}
%     \hfill\qedsymbol
%     \end{minipage}
%     \end{trivlist}
% }

\newcommand*\circled[1]{\tikz[baseline=(char.base)]{
		\node[shape=circle,draw,inner sep=1pt] (char) {#1};}}
\newcommand\getcurrentref[1]{%
	\ifnumequal{\value{#1}}{0}
	{??}
	{\the\value{#1}}%
}

\newcommand{\getCurrentSectionNumber}{\getcurrentref{section}}
\newenvironment{myproof}[1][\proofname]{%
	\proof[\bfseries #1: ]%
}{\endproof}

\newcounter{mylabelcounter}

\makeatletter
\newcommand{\setword}[2]{%
	\phantomsection
	#1\def\@currentlabel{\unexpanded{#1}}\label{#2}%
}
\makeatother

%%%%%%%%%%%%%%%%%%%%%%%%%%%%%%%%%%%%%%%%%%%
% TABLE OF CONTENTS
%%%%%%%%%%%%%%%%%%%%%%%%%%%%%%%%%%%%%%%%%%%

\usepackage{tikz}
\definecolor{doc}{RGB}{0,60,110}
\usepackage{titletoc}
\contentsmargin{0cm}
\titlecontents{chapter}[3.7pc]
{\addvspace{30pt}%
	\begin{tikzpicture}[remember picture, overlay]%
		\draw[fill=doc!60,draw=doc!60] (-7,-.1) rectangle (-0.9,.5);%
		\pgftext[left,x=-3.5cm,y=0.2cm]{\color{white}\Large\sc\bfseries Chapter\ \thecontentslabel};%
	\end{tikzpicture}\color{doc!60}\large\sc\bfseries}%
{}
{}
{\;\titlerule\;\large\sc\bfseries Page \thecontentspage
	\begin{tikzpicture}[remember picture, overlay]
		\draw[fill=doc!60,draw=doc!60] (2pt,0) rectangle (4,0.1pt);
	\end{tikzpicture}}%
\titlecontents{section}[3.7pc]
{\addvspace{2pt}}
{\contentslabel[\thecontentslabel]{2pc}}
{}
{\hfill\small \thecontentspage}
[]
\titlecontents*{subsection}[3.7pc]
{\addvspace{-1pt}\small}
{}
{}
{\ --- \small\thecontentspage}
[ \textbullet\ ][]

\makeatletter
\renewcommand{\tableofcontents}{%
	\chapter*{%
	  \vspace*{-20\p@}%
	  \begin{tikzpicture}[remember picture, overlay]%
		  \pgftext[right,x=15cm,y=0.2cm]{\color{doc!60}\Huge\sc\bfseries \contentsname};%
		  \draw[fill=doc!60,draw=doc!60] (13,-.75) rectangle (20,1);%
		  \clip (13,-.75) rectangle (20,1);
		  \pgftext[right,x=15cm,y=0.2cm]{\color{white}\Huge\sc\bfseries \contentsname};%
	  \end{tikzpicture}}%
	\@starttoc{toc}}
\makeatother


